\subsection{\texorpdfstring{半范数 \(\mathbf{N}_p\)}{半范数 N sub p}}

在以下全部内容中，F 表示域 R 或域 C 上的一个完备赋范向量空间（即 Banach 空间）；元素 \(z \in F\) 的范数将记作 \(|z|\) 。给定集合 A 到 F 中的一个映射 f，我们用 \(|f|\) 表示映射 \(x \mapsto |\mathbf{f}(x)|\) ，即 A 到 \(R_{+}\) 中的一个映射（必须注意 \(|f|\) 是一个数值函数，而不是一个数）。

定义 1. --- 设 X 是一个局部紧空间，\(\mu\) 是 X 上的一个测度。对 X 到 Banach 空间 F 中的每个映射 f，以及对每个使 \(1 \leqslant p < +\infty\) 的数 p，我们用 \(\mathrm{N}_p(\mathbf{f}, \mu)\) （或简作 \(\mathrm{N}_p(\mathbf{f})\) ）表示正数 \((\int^* |\mathbf{f}|^p d|\mu|)^{1/p}\) 。

注意，数 \(N_{p}(\mathbf{f})\) 可以等于 \(+\infty\) 。

命题 2. --- 若 f 与 g 是 X 到 F 中的两个映射，且 \(\alpha\) 是任一标量 \(\neq 0\) ，则对 \(1 \leqslant p < +\infty\) ，

\begin{enumerate}
\def\labelenumi{(\arabic{enumi})}
\setcounter{enumi}{2}
\tightlist
\item
\end{enumerate}

\[
\mathrm{N} _ {p} (\alpha \mathbf {f}) = | \alpha | \mathrm{N} _ {p} (\mathbf {f})\tag{4}
\]

\[
\mathrm{N} _ {p} (\mathbf {f} + \mathbf {g}) \leqslant \mathrm{N} _ {p} (\mathbf {f}) + \mathrm{N} _ {p} (\mathbf {g}).
\]

因为，关系式 (3) 由定义 1 以及 \(|\mu|^*\) 是正齐次的这一事实立得；另一方面，由于 \(|\mathbf{f} + \mathbf{g}| \leqslant |\mathbf{f}| + |\mathbf{g}|\) ，不等式 (4) 由 Minkowski 不等式 (1) 以及 \(|\mu|^*\) 是递增的这一事实得出。

我们把定义 1 推广到定义在 X 上的数值函数（有限或非有限）的情形，仍令

\[
\mathrm{N} _ {p} (f) = \left(\int^ {*} | f | ^ {p} d | \mu |\right) ^ {1 / p}
\]

对这样的函数 \(f\) 。立即可见，当 \(f + g\) 在 \(X\) 上有定义且 \(\alpha \neq 0\) 时，关系式 (3) 与 (4) 对这些函数也成立。此外：

定理 1（可数凸性定理）. --- 设 \((f_n)\) 是定义在 X 上的函数 \(\geqslant 0\) （有限或非有限）的一个序列。对 \(1 \leqslant p < +\infty\) ，

\[
\mathrm{N} _ {p} \left(\sum_ {n = 1} ^ {\infty} f _ {n}\right) \leqslant \sum_ {n = 1} ^ {\infty} \mathrm{N} _ {p} (f _ {n}).\tag{5}
\]

令 \(f = \sum_{n=1}^{\infty} f_n\) ；\(f\) 是递增函数序列 \(g_n = \sum_{k=1}^{n} f_k\) 的上包络；\(N sub p(f)\) 的定义以及 §1, No.~3 的定理 3 表明 \(\mathrm{N}_{p}(f)=\sup_{n}\mathrm{N}_{p}(g_{n})\) 。但由命题 1 有 \(\mathrm{N}_{p}(g_{n})\leqslant\sum_{k=1}^{n}\mathrm{N}_{p}(f_{k})\) ，由此得不等式 (5)。

命题 3. --- 若 \(\mathbf{f}\) 与 \(\mathbf{g}\) 是 X 到 Banach 空间 F 中的两个等价映射，则 \(\mathrm{N}_p(\mathbf{f} - \mathbf{g}) = 0\) 对 \(1 \leqslant p < +\infty\) 成立；反之，若 \(\mathrm{N}_p(\mathbf{f} - \mathbf{g}) = 0\) 对某个值 \(p \geqslant 1\) 成立，则 \(\mathbf{f}\) 与 \(\mathbf{g}\) 等价。命题由 §2, No.~3 的定理 1 立得。

若 \(\mathbf{f}\) 与 \(\mathbf{g}\) 是 \(X\) 到 \(F\) 中的两个等价映射，则 \(\mathrm{N}_p(\mathbf{f}) = \mathrm{N}_p(\mathbf{g})\) 对一切 \(p \geqslant 1\) 成立（ \(\S 2\) , No.~3, 命题 6）；于是 \(\mathrm{N}_p(\mathbf{f})\) 只依赖于类 \(\widetilde{\mathbf{f}}\) （\(\mathbf{f}\) 的），且依定义令 \(\mathrm{N}_p(\widetilde{\mathbf{f}}) = \mathrm{N}_p(\mathbf{f})\) 。由于 \(X\) 到 \(F\) 中的映射的类构成一个向量空间（ \(\S 2\) , No.~4），关系式 (3) 与 (4) 也可写成

\begin{enumerate}
\def\labelenumi{(\arabic{enumi})}
\setcounter{enumi}{5}
\tightlist
\item
\end{enumerate}

\[
\mathrm{N} _ {p} (\alpha \widetilde {\mathbf {f}}) = | \alpha | \mathrm{N} _ {p} (\widetilde {\mathbf {f}})\tag{7}
\]

\[
\mathrm{N} _ {p} (\widetilde {\mathbf {f}} + \widetilde {\mathbf {g}}) \leqslant \mathrm{N} _ {p} (\widetilde {\mathbf {f}}) + \mathrm{N} _ {p} (\widetilde {\mathbf {g}}).
\]

类似地，对等价数值函数（有限或非有限）的每个类，定义 \(\mathrm{N}_{p}(\tilde{\mathbf{f}})\) 。

于是，可对下述函数定义 \(\mathrm{N}_{p}(\mathbf{f})\) ：该函数在 X 上几乎处处有定义，取值于 F （相应地 \(\overline{R}\) ）中，方法是令 \(\mathrm{N}_{p}(\mathbf{f})=\mathrm{N}_{p}(\widetilde{\mathbf{f}})\) ；此时显然关系式 (3) 与 (4) 仍成立（在数值函数（有限或非有限）的情形，假定 \(\alpha\neq0\) 且 \(f+g\) 几乎处处有定义）。

若 \(0 < p < 1\) ，仍令 \(\mathrm{N}_p(\mathbf{f}) = (\int^* |\mathbf{f}|^p d|\mu|)^{1/p}\) ，但不等式 (4) 与 (5) 不再成立（参见 Ch. I, 习题 6 与 Ch. IV, §6, 习题 13）。

